\documentclass{article}
\usepackage[utf8]{inputenc}
\usepackage[spanish]{babel}
\usepackage{listings}
\usepackage{xcolor}
\usepackage{geometry}
\usepackage{amsmath}
\geometry{a4paper, margin=1in}
\definecolor{codegreen}{rgb}{0,0.6,0}
\definecolor{codegray}{rgb}{0.5,0.5,0.5}
\definecolor{codepurple}{rgb}{0.58,0,0.82}
\definecolor{backcolour}{rgb}{0.95,0.95,0.92}
\lstset{backgroundcolor=\color{backcolour},commentstyle=\color{codegreen},keywordstyle=\color{magenta},numberstyle=\tiny\color{codegray},stringstyle=\color{codepurple},basicstyle=\ttfamily\small,breakatwhitespace=false,breaklines=true,captionpos=b,keepspaces=true,numbers=left,numbersep=5pt,showspaces=false,showstringspaces=false,showtabs=false,tabsize=2,language=C++}
\newcommand{\quees}[1]{\section*{🤔 ¿QUÉ ES?}\hrulefill \\ #1}
\newcommand{\dichodeotraforma}[1]{\section*{💬 DICHO DE OTRA FORMA}\hrulefill \\ #1}
\newcommand{\diagrama}[1]{\section*{🎨 DIAGRAMA}\hrulefill \\ #1}
\newcommand{\comofunciona}[1]{\section*{⚙️ ¿CÓMO FUNCIONA?}\hrulefill \\ #1}
\newcommand{\complejidad}[1]{\section*{⏱️ COMPLEJIDAD (Big-O)}\hrulefill \\ #1}
\newcommand{\entrevistas}[1]{\section*{🎯 EN ENTREVISTAS}\hrulefill \\ #1}
\newcommand{\resumen}[1]{\section*{📋 RESUMEN RÁPIDO}\hrulefill \\ #1}
\begin{document}
\title{Introducción a Dynamic Programming (DP)}
\date{}
\maketitle
\quees{Un estudiante inteligente que guarda sus apuntes 📒. 

Imagina que en un examen te preguntan 50 veces cuánto es 123 x 456.
En lugar de hacer la multiplicación cada vez (y perder tiempo),
la haces la primera vez, anotas el resultado en tu cuaderno, y 
las siguientes 49 veces ¡solo lo consultas!

Dynamic Programming es simplemente eso: guardar resultados de 
sub-problemas ya resueltos para no repetir trabajo innecesariamente. 
Es el arte de intercambiar un poco de memoria por mucho tiempo.}
\dichodeotraforma{Dynamic Programming (DP) es una técnica de optimización para resolver
problemas complejos dividiéndolos en sub-problemas más pequeños. 

Para usar DP, el problema debe cumplir dos condiciones:
1. Overlapping Subproblems: El mismo sub-problema se calcula varias
   veces (como en Fibonacci).
2. Optimal Substructure: La solución óptima del problema global se
   puede construir a partir de las soluciones óptimas de sus partes.}
\diagrama{\begin{verbatim}
Sin DP (Calculas lo mismo muchas veces):
             fib(5)
           /        \
      fib(4)         fib(3)
      /    \        /      \
  fib(3)  fib(2)  fib(2)  fib(1)
  /    \
fib(2) fib(1)

¡Nota cómo fib(3) y fib(2) se repiten! 💀

Con DP (Memoization):
Calculas fib(3) una vez, lo guardas, y cuando fib(4) lo necesite,
simplemente lo lee de la memoria. ⚡
\end{verbatim}}
\comofunciona{Existen dos formas de implementar DP:

1. Top-Down (Memoization):
   - Empiezas con el problema grande.
   - Lo divides recursivamente.
   - Guardas el resultado en un Hash Map o Array antes de regresarlo.

2. Bottom-Up (Tabulation):
   - Empiezas con los casos más pequeños (base cases).
   - Llenas una tabla (array) paso a paso hasta llegar al resultado.
   - Generalmente se hace con un loop iterativo.}
\section*{📝 PSEUDOCÓDIGO}
\hrulefill
\begin{verbatim}
// Fibonacci Sin DP: O(2^n) - Pésimo
función fib(n):
    SI n <= 1: return n
    return fib(n-1) + fib(n-2)

// Fibonacci Con DP (Memo): O(n) - Excelente
función fib(n, memo):
    SI n EN memo: return memo[n]
    SI n <= 1: return n
    memo[n] = fib(n-1, memo) + fib(n-2, memo)
    return memo[n]
\end{verbatim}
\section*{💻 CÓDIGO C++}
\hrulefill
\begin{lstlisting}[language=C++]
#include <iostream>
#include <vector>
using namespace std;

// Fibonacci sin DP (muy lento para n > 40)
long long fibNaive(int n) {
    if (n <= 1) return n;
    return fibNaive(n - 1) + fibNaive(n - 2);
}

// Fibonacci con DP (Tabulation - Bottom Up)
long long fibDP(int n) {
    if (n <= 1) return n;
    vector<long long> dp(n + 1);
    dp[0] = 0;
    dp[1] = 1;
    for (int i = 2; i <= n; i++) {
        dp[i] = dp[i - 1] + dp[i - 2];
    }
    return dp[n];
}

int main() {
    int n = 50;
    cout << "Fibonacci de " << n << " es: " << fibDP(n) << endl;
    return 0;
}
\end{lstlisting}
\complejidad{Enfoque        │ Tiempo    │ Espacio
  ───────────────┼───────────┼────────────────────────────────
  Fuerza Bruta   │ O(2\textasciicircum{}n)    │ O(n) (stack)
  DP (Memo/Tab)  │ O(n)      │ O(n) (tabla)

* La DP reduce la complejidad de exponencial a lineal en muchos casos.}
\entrevistas{}
\resumen{• DP = Recursión + Memoria.
• No calcules lo mismo dos veces.
• Overlapping Subproblems: La clave para usar DP.
• Top-Down: Fácil de pensar (recursivo).
• Bottom-Up: Más eficiente (iterativo).
• Esencial para problemas de optimización.
════════════════════════════════════════════════════════════════}
\end{document}